\section{The highest weight theorem}

Let $G$ be a semisimple algebraic group, $T$ a maximal torus, $B = TU$ a Borel subgroup containing $T$, where $U = B_u$. Since, in characteristic $0$, every representation of $G$ is made up of irreducible representations, we wish to understand these precisely. This chapter aims to give the required tools to tackle the highest weight theorem, essentially classifying all possible irreducible representations of $G$.

\subsection{Weights of a representation}
Let $G$ be a semisimple algebraic group $G$, $T$ a maximal torus of $G$, $B = TU$ a Borel subgroup containing $T$.

Morphisms into the general linear group of a vector space have a specific name, as is the case in classical group theory.
\begin{defn}
    A morphism of algebraic groups $\rho: G \to \GL{}{V}$ is called a \emph{(rational) representation} of $G$, or \emph{$G$-module}. It is \emph{irreducible} if it has no non-trivial $G$-stable subspace or, equivalently, if it is not the direct sum of two $G$-modules.
\end{defn}
Let $\rho: G \to \GL{}{V}$ be a representation of $G$, with $V$ a finite-dimensional vector space. Since $\rho(T)$ is a torus in $\GL{}{V}$, we may look at its weights. By definition a weight of $\rho(T)$ is a map $\alpha: \rho(T) \to \Gm$ such that $V_\alpha = \left\{v \in V ~\middle|~ t v = \alpha(t) v \textnormal{ for all } t \in \rho(T) \right\}$ is non-trivial. By precomposition with $\rho$ any character of $\rho(T)$ gives a character of $T$. This is in particular the case for weights of $\rho(T)$, motivating the following definition.
\begin{defn}
    Let $\rho: G \to \GL{}{V}$ be a representation of $G$, with $V$ a finite-dimensional vector space. The \emph{weights} of $\rho$ are the images in $X(T)$ of the weights of $\rho(T)$ by the precomposition $X(\rho(T)) \xrightarrow{- \circ \rho} X(T)$.
\end{defn}

Since a map $\rho : G \to \GL{}{V}$ is precisely a $G$-action on $V$, a weight space $V_\lambda$ may be simply described as $\left\{v \in V ~\middle|~ t.v = \lambda(t) v \text{ for all } t \in T\right\}$, for a weight $\lambda \in X(T)$. The \emph{multiplicity} of a weight is $\dim V_\lambda$.

\begin{exmp}[The exterior power representation of $\SL{n}{K}$]\label{ex:sym_rep}
    Let ${\rho: \SL{n}{K} \to \GL{}{\wedge^k K^n}}$ be the representation induced by the $\SL{n}{K}$-action on exterior powers (\cref{subsubsection:actions}). Let $t = \textnormal{diag}(t_1, \ldots, t_n)\in T$ be a diagonal matrix in $\SL{n}{K}$. Studying its action on a basis of $\wedge^k K^n$ will allow us to find the weights of $\rho$. A basis of $\wedge^k K^n$ is given by products $e_{i_1} \wedge \cdots \wedge e_{i_k}$ where $i_1 < i_2 < \cdots < i_k$. The action of $t$ is then
    \[ t . e_{i_1} \wedge \cdots \wedge e_{i_k} = (t e_{i_1}) \wedge \cdots \wedge (t e_{i_k}) = (t_{i_1} e_{i_1}) \wedge \cdots \wedge (t_{i_k} e_{i_k}) = \prod_{j = 1}^k \varepsilon_{i_j}(t) \cdot (e_{i_1} \wedge \cdots \wedge e_{i_k}).\]
    Thus the weights of $\rho$ are sums $\sum_{j = 1}^k \varepsilon_{i_j}$ where each $i_j$ is distinct.
\end{exmp}
\subsection{One-parameter subgroups}

In order to define positive roots, we will need a way to associate an integer to a root, given a specific Borel subgroup. To do so, the notion of one-parameter subgroup is of paramount importance. 

\subsubsection{Pairing one-parameter subgroups and weights}

Just as characters are an essential tool for the study of algebraic groups, the dual notion is also of particular interest.
\begin{defn}
    Let $T$ be a torus. A \emph{one-parameter (multiplicative) subgroup} of $T$ is a morphism of algebraic groups $\lambda : \Gm \rightarrow T$. The set of one-parameter subgroups is denoted $Y(T)$.
\end{defn}
Just like characters inherited the structure of an abelian group, one-parameter subgroups also get one by $(\lambda \mu)(a) = \lambda(a)\mu(a)$. As we have defined maps to and from $\Gm$, we may compose them to get a map $\Gm \to \Gm$, which is exactly a character of $\Gm$. Since ${X(\Gm) \cong \ZZ}$, this yields a pairing 
\begin{alignat*}{4}
    \langle \cdot , \cdot \rangle: \quad & X(T) \times Y(T) &\ \to\ &\ZZ\\
           & \quad (\chi, \lambda) &\ \mapsto\ &\langle \chi, \lambda \rangle = \chi \circ \lambda.
\end{alignat*}

\begin{exmp}[One-parameter subgroups of the standard maximal torus of $\SL{n}{K}$]\label{ex:1-psg}
    Let $G = \SL{n}{K}$ and $T$ be the torus of diagonal matrices with determinant $1$. A one-parameter subgroup is a map $\lambda: \Gm \to T$ that is a morphism of varieties. The definition of morphism of varieties requires that $\lambda$ is of the type $\lambda(t) = \textnormal{diag}(t^{i_1}, \ldots, t^{i_n})$ for some integers $i_k$. The determinant of the matrix being $1$ requires $\sum i_j = 0$. As in the case of $X(T)$, this means that $Y(T)$ is free abelian on $n-1$ generators.\\
    To compute the pairing, fix a $\lambda \in Y(T)$, i.e.~a map $t \mapsto \textnormal{diag}(t^{i_1}, \ldots, t^{i_n})$ with $\sum i_j = 0$. Let $\chi \in X(T)$ be a character of $T$, i.e.~a sum $\sum n_i \cdot \varepsilon_i$ such that $\sum n_i = 0$. The pairing of $\lambda$ with $\chi$ is the composite
    \[ t \mapsto \textnormal{diag}(t^{i_1}, \ldots, t^{i_n}) \mapsto \left(\sum n_i \cdot \varepsilon_i\right) (\textnormal{diag}(t^{i_1}, \ldots, t^{i_n})) = \prod_{j = 1}^{n}{\left(t^{i_j}\right)}^{n_j} = t^{\sum_{j=1}^{n} i_j \cdot n_j}.\]
    In particular, the identification $X(\Gm) \cong \ZZ$ says that $\langle \chi, \lambda\rangle = \sum_{j=1}^{n} i_j \cdot n_j$.
\end{exmp}

\subsubsection[The action of a one-parameter subgroup]{The action of a one-parameter subgroup~\cite[Section~25.1]{humphreys_linear_1975}}

For this part, let $T$ be a torus in $\GL{}{V}$ with a one-parameter subgroup $\lambda: \Gm \to T$. 
This one-parameter subgroup yields an action of $\Gm$ on $\PP\left(V\right)$ as follows. 
Let $(v_1, \ldots, v_n)$ be a basis of $V$ consisting of eigenvectors for $T$. 
Consider the weights of $T$ corresponding to each of these vector $\varepsilon_1, \ldots, \varepsilon_n$. 
These weights generate $X(T)$. 
Let $m_i = \langle \varepsilon_i, \lambda \rangle$. 
For a vector $v = \sum_{i = 1}^{n} a_i v_i \in V$ and an element $t \in \Gm$, we have by definition
\[ \lambda(t).v = \sum_{i = 1}^n a_i \lambda(t).v_i = \sum_{i = 1}^n a_i \varepsilon_i ( \lambda(t)) v_i = \sum_{i = 1}^n a_i \langle \varepsilon_i, t \rangle v_i = \sum_{i = 1}^n a_i t^{m_i} v_i \]
where, in the last equality, we use explicitly the isomorphism $\ZZ \cong X(\Gm)$. We therefore define the action of $\Gm$ on $\PP\left(V\right)$ by $t . [v] = [\lambda(t). v]$.

For a fixed non-zero vector $v \in V$, let $\phi: \Gm \to \PP\left(V\right)$ be defined as $\phi(t) = t. [v]$. This map can be extended to a morphism $\proj{1}{K} \to \PP\left(V\right)$, which we view as a gluing of $K^* \cup \{0\}$ and $K^* \cup \{\infty\}$. These correspond to the standard affine opens $D_+(x_0)$ and $D_+(x_1)$. By identifying $\Gm$ with $K^*$, we need only define two maps $\phi_0 : K^* \cup \{0\} \to \PP\left(V\right)$ and $\phi^0 : K^* \cup \{\infty\} \to \PP\left(V\right)$ that agree with $\lambda$ applied to $v$ on $K^*$. \\
To do so, let $m_0 = \min{\left\{m_i ~\middle|~ i \in I\right\}}$ and $m^0 = \max{\left\{m_i ~\middle|~ i \in I\right\}}$, where $I = \left\{i ~\middle|~ a_i \neq 0\right\}$. Define $I_0 = \left\{ i \in I ~\middle|~ m_i = m_0\right\}$ and $I^0 = \left\{ i \in I ~\middle|~ m_i = m^0\right\}$.\\
Intuitively, $\phi_0(0)$ should be obtained as a limit in $\Gm$ tending to zero, while $\phi^0(\infty)$ should be a limit tending to infinity. Asymptotic behavior would then suggest that all indices not in $I_0$ are ignored while taking the limit to zero (since $t^{m_i}$ would then tend to zero far quicker), while all indices not in $I^0$ are ignored when tending to infinity. This suggests the values $\phi_0(0) = [\sum_{i \in I_0} a_i v_i]$ and ${\phi^0(\infty) = [\sum_{i \in I^0} a_i v_i]}$, where projectification discards the scalars $t^{m_i}$ (since the $m_i$ are all the same when in $I_0$ or $I^0$). \\
Formally, the maps 
\[\phi_0(t) := \left[\sum_{i=0}^{n} a_i t^{m_i - m_0} v_i\right] \quad \text{and} \quad \phi^0(t) := \left[\sum_{i=0}^{n} a_i t^{m_i - m^0} v_i\right]\]
are well-defined on $K^* \cup \{0\}$ and $K^* \cup \{\infty\}$ respectively (since all exponents are positive, respectively negative). Furthermore, they coincide with $\phi$ on $K^*$, since the scalar $t^{-m_0}$ (respectively $t^{-m^0}$) is taken care of via projectification. Glueing them therefore yields the desired map $\PP^1_K \to \PP\left(V\right)$. For easier understanding, the images of $0$ and $\infty$ will be denoted $\lambda(0)[v]$ and $\lambda(\infty)[v]$. Their explicit formulae are
\[ \lambda(0)[v] = \left[\sum_{i \in I_0} a_i v_i\right] \quad \text{and} \quad \lambda(\infty)[v] = \left[\sum_{i \in I^0} a_i v_i\right].\]

Note that $t \cdot \lambda(0)[v] = \left[\sum_{i \in I_0} a_i t^{m_0} v_i\right] = \left[\sum_{i \in I_0} a_i v_i\right]$ for any $t \in \Gm$ (and similarly for $\lambda(\infty)[v]$) for any $t \in \lambda(\Gm)$, so that $\lambda(\Gm)$ fixes these points.

\subsection{Positive roots}
\subsubsection{Regular tori and one-parameter subgroups}

Recall that every torus is contained in at least one Borel subgroup. One may show in particular that maximal tori are contained in a finite number of Borel subgroups. This motivates the following definition.
\begin{defn}
    Let $T$ be a torus in an algebraic group $G$. Let $\mathfrak{B}^T$ denote the set of Borel subgroups containing $T$. The torus $T$ is called \emph{regular} if $\mathfrak{B}^T$ is finite. Else, it is \emph{singular}.
\end{defn}
Thus, maximal tori are always regular.

The regularity of a torus is completely determined by the solvability of its centralizer, as is proven in~\cite[Section~24.2]{humphreys_linear_1975}.
\begin{theorem}\label{thm:regular_torus}
    Let $S$ be a torus in $G$. Then $S$ is regular if and only if $C_G(S)$ is solvable.
\end{theorem}

Let $\lambda: \Gm \to T$ be a one-parameter subgroup. Since $\Gm$ is trivially a torus, the image $\lambda(\Gm)$ is also a torus, motivating the following definition.
\begin{defn}
    Let $\lambda \in Y(T)$ be a one-parameter subgroup. If $\lambda(\Gm) \subseteq G$ is a regular torus, the one-parameter subgroup $\lambda$ is called \emph{regular}. Else, it is \emph{singular}. \\
    The set of all regular one-parameter subgroups is denoted ${Y(T)}_{\textnormal{reg}}$.
\end{defn}

\begin{exmp}[Regular one-parameter subgroups in $\SL{n}{K}$]\label{ex:reg-1psg}
    We continue the computations from \cref{ex:1-psg}, where we showed that a one-parameter subgroup of the maximal torus of diagonal matrices is a map $t \mapsto \textnormal{diag}(t^{l_1}, \ldots, t^{l_n})$ with $\sum k_i = 0$. \\
    Let us use \cref{thm:regular_torus} to find which one-parameter subgroups are regular. Note that $C_G(\lambda(\Gm))$ is the intersection of the centralizer of each $\lambda(t)$. \\
    Fix a matrix $\lambda(t) = \textnormal{diag}(t^{l_1}, \ldots, t^{l_n})$, and let $A = (a_{ij}) \in \SL{n}{K}$. Note that 
    \[ {(\lambda(t) \cdot A )}_{ij} = t^{l_i} a_{ij} \]
    while 
    \[ {(A \cdot \lambda(t))}_{ij} = t^{l_j} a_{ij}. \]
    Thus a matrix $A$ commutes with $\lambda(t)$ if and only if
    \[ t^{l_j} a_{ij} = t^{l_i} a_{ij} \quad \iff \quad a_{ij} (t^{l_j} - t^{l_i}) = 0 \]
    for all $i, j \in \{1, \ldots, n \}$. In particular, a matrix $A$ is in $C_G(\lambda(\Gm))$ if and only if, for all $i,j$, either $l_i = l_j$ or $a_{ij} = 0$. In particular, if there are distinct $i,j$ such that $l_i = l_j$, then a copy of $\SL{2}{K}$ embeds into $C_G(\lambda(\Gm))$ (into coordinates $(i,i)$, $(i, j)$, $(j,i)$ and $(j,j)$), so that $C_G(\lambda(\Gm))$ is not solvable, and $\lambda$ is not regular. On the other hand, if all $l_i$'s are distinct, then $C_G(\lambda(\Gm)) = T$, which is abelian, so solvable. \\
    Thus, a one-parameter subgroup $\lambda \in Y(T)$ parametrized by integers $(l_1, \ldots, l_n)$ that sum to zero is regular if and only if all integers $l_i$ are distinct.
\end{exmp}

\subsubsection{The Borel subgroup of a regular one-parameter subgroup}

Following the discussion in~\cite[Section~25.4]{humphreys_linear_1975}, the action of $\lambda \in {Y(T)}_{\textnormal{reg}}$ on $\PP(V)$ constructed above allows one to attach to $\lambda$ a canonical Borel subgroup $B(\lambda) \in \mathfrak{B}^T$.

Choose an embedding $G \hookrightarrow \GL{}{V}$ such that $G/B$ is identified with the $G$-orbit $X$ of some point of $\PP(V)$, with $X$ contained in no hyperplane $\PP(W)$ (this can be done by \cref{thm:chevalley}). For $\lambda \in {Y(T)}_{\textnormal{reg}}$, let $(\gamma_1, \ldots, \gamma_n)$ be the weights of $T$ on $V$ and $(v_1, \ldots, v_n)$ an ordered eigenbasis such that
\[ m^0 = m_1 = \cdots = m_r > m_{r+1} \geq \cdots \geq m_n = m_0, \quad \text{where } m_i = \langle \gamma_i, \lambda \rangle. \]
Since $X$ lies in no hyperplane, we may pick $[v] = \bigl[\sum_i a_i v_i\bigr] \in X$ with $a_1 \neq 0$, so the formula of the preceding section gives $\lambda(\infty)[v] = [a_1 v_1 + \cdots + a_r v_r]$.

\begin{lemma}
    If $\lambda$ is regular, then $r = 1$.
\end{lemma}
\begin{proof}
    If $r > 1$, then for any distinct $b, b' \in K$, the points $[v_1 + bv_2 + \cdots]$ and $[v_1 + b' v_2 + \cdots]$ both lie in $X$ (since $X$ avoids only finitely many hyperplanes of the form ``ratio of second to first coordinate equals $c$\,''), and both are fixed by $\lambda(\Gm)$. This yields infinitely many distinct fixed points of $\lambda(\Gm)$ in $X$, contradicting the regularity of $\lambda(\Gm)$, since points of $X$ are in bijection with Borel subgroups of $G$.
\end{proof}

Therefore $\lambda(\infty)[v] = [v_1]$ is the unique point of $X$ fixed by $\lambda(\Gm)$. Since $G/B$ corresponds to $X$, this defines a unique Borel subgroup of $G$.

\begin{defn}
    For $\lambda \in {Y(T)}_{\textnormal{reg}}$, the \emph{Borel subgroup associated to $\lambda$} is
    \[ B(\lambda) = \operatorname{Stab}_G\bigl(\lambda(\infty)[v]\bigr) = \operatorname{Stab}_G([v_1]). \]
\end{defn}
Since $T$ fixes $[v_1]$ (as $\lambda$ is regular), we have $B(\lambda) \in \mathfrak{B}^T$. Furthermore, we may show that any Borel subgroup $B \in \mathfrak{B}^T$ can be obtained this way. 

One can thus define the opposite Borel subgroup. If $B = B(\lambda)$, its opposite Borel subgroup is the group $B^{-} = B(-\lambda)$.

\begin{exmp}[The regular one-parameter subgroup of the standard Borel subgroup of $\SL{n}{K}$]\label{ex:reg_1psg_borel}
    Let $G = \SL{n}{K}$, $T$ be the maximal torus of diagonal matrices and $B$ the standard Borel subgroup of upper triangular matrices. In \cref{ex:reg-1psg}, it was established that a one-parameter subgroup $\lambda = (l_1, \ldots, l_n)$ is regular if and only if each $l_i$ is different. Let us find a regular one-parameter subgroup such that $B = B(\lambda)$.

    To do so, recall the variety $G/B$ may be identified with the flag variety $\flag{K^n}$, which is closed in some projective space $\PP(V)$, and is the $G$-orbit of the standard flag ${0 \subseteq \vspan{e_1} \subseteq \cdots \subseteq \vspan{e_1, \ldots, e_n}}$. $V$ can in particular be chosen to be ${\wedge^1 K^n \otimes \cdots \otimes \wedge^{n-1} K^n}$, so that $G$ acts simultaneously on each exterior power as described in \cref{subsubsection:actions}. In that context, $G/B$ is seen as the orbit of the element
    \[ v := e_1 \oplus \left(e_1 \wedge e_2 \right) \oplus \cdots \oplus \left(e_1 \wedge \cdots \wedge e_{n-1}\right)\]
    seen in the projectification. In particular, finding a one-parameter subgroup $\lambda$ such that the pairing $\langle \varepsilon_i, \lambda \rangle$ is strictly decreasing with $i$ is sufficient. One may choose for example $\lambda(t) = \textnormal{diag}(t^{n-1}, t^{n-3},\ldots, t^{-(n-1)})$ if $n$ is even, and $\lambda(t) = \textnormal{diag}(t^{n}, t^{n-2},\ldots, t^{-(n)})$ if $n$ is odd. Doing so, the action of $\lambda(\infty)$ on $v$ fixes $e_1$ in the first component, fixes $e_1 \wedge e_2$ in the second, etc.. Thus $\lambda(\infty)([v]) = [v]$, and $B(\lambda) = \stab{G}{[v]} = B$.
\end{exmp}

\subsubsection{Positive roots with respect to a Borel subgroup}

Let $B$ be a Borel subgroup of $G$ containing a maximal torus $T$. By the previous section, there is a regular one-parameter subgroup $\lambda \in {Y(T)}_{reg}$ such that $B = B(\lambda)$. The pairing $\langle \cdot , \cdot \rangle: X(T) \times Y(T) \to \ZZ$ gives a way to decide when a root $\alpha \in \Phi$ is positive by fixing the second parameter to $\lambda$.

\begin{defn}
    Let $T$ be a maximal torus of $G$, $B = B(\lambda)$ be the Borel subgroup corresponding to $\lambda \in {Y(T)}_{reg}$. A root $\alpha \in \Phi = \Phi(G,T)$ is called \emph{positive} if $\langle \alpha, \lambda \rangle > 0$. Else it is called \emph{negative}. The set of positive roots is denoted $\Phi^+$, while the set of negative roots is denoted $\Phi^{-}$.
\end{defn}

Among the positive roots, those which cannot be expressed as a sum of two positive roots are called \emph{simple roots}. They form the \emph{base} $\Delta = \{\alpha_1, \ldots, \alpha_l\}$ of $\Phi$. Every root in $\Phi^+$ (resp.\ $\Phi^-$) is a non-negative (resp.\ non-positive) integer linear combination of simple roots.

\begin{exmp}[Positive and simple roots of $\SL{n}{K}$]
    For ease of calculations, let $n$ be even Let ${\lambda \in {Y(T)}_{reg}}$ be the one-parameter subgroup given by $t \mapsto \textnormal{diag}(t^{n-1}, t^{n-3}, \ldots, t^{-(n-1)})$ (see~\cref{ex:reg_1psg_borel}). Recall the roots of $\SL{n}{K}$ are given by $\varepsilon_i - \varepsilon_j$, $i \neq j$. The computations in \cref{ex:1-psg} yield that
    \[ \langle \varepsilon_i - \varepsilon_j, \lambda \rangle = (n + 1 - 2i) \cdot 1 + (n + 1 - 2j) \cdot (-1) = 2j - 2i.\]
    In particular, the result is positive if and only if $j > i$, giving a partition of $\Phi(\SL{n}{K}, T)$ into $\Phi^{+} = \left\{\epsilon_i - \epsilon_j ~\middle|~ 1 \leq i < j \leq n \right\}$ and $\Phi^{-} = \left\{\epsilon_i - \epsilon_j ~\middle|~ 1 \leq j < i \leq n\right\}$. The basis of $\Phi^{+}$ is given by the roots $\Delta = \left\{ \epsilon_i - \epsilon_{i+1} ~\middle|~ 1 \leq i \leq n-1 \right\}$.
\end{exmp}


\subsection{Highest weight theorem}

\subsubsection{Dominant weights}

With $\Phi^+$ and $\Delta = \{\alpha_1, \ldots, \alpha_\ell\}$ determined by $B$ as above, we identify the class of weights that will parametrise irreducible representations of $G$.

Recall from \cref{thm:root_system} that $X(T)$ embeds as a full-rank lattice in
$E = X(T) \otimes_\ZZ \RR$, and that each root $\alpha \in \Phi$ yields Cartan integers
$\langle \chi, \alpha^\vee \rangle = \frac{2(\chi, \alpha)}{(\alpha, \alpha)} \in \ZZ$
for all $\chi \in X(T)$, where $(\cdot, \cdot)$ is the $W$-invariant inner product on $E$.

\begin{defn}
    A weight $\chi \in X(T)$ is \emph{dominant} with respect to $B$ if
    $\langle \chi, \alpha_i^\vee \rangle \geq 0$ for every simple root $\alpha_i \in \Delta$. The set of all dominant weights is denoted ${X(T)}_+$.
\end{defn}

Since each element of $\Phi^+$ is a non-negative integer combination of simple roots
and $\langle \chi, \cdot \rangle$ is $\ZZ$-linear, the condition on $\Delta$ implies
$\langle \chi, \alpha^\vee \rangle \geq 0$ for all $\alpha \in \Phi^+$; it therefore
suffices to check dominance on $\Delta$ alone.

There is a canonical basis of $E$ dual to the simple coroots. The
\emph{fundamental dominant weights} $\omega_1, \ldots, \omega_\ell \in E$ are the
unique elements satisfying
\[
    \langle \omega_i, \alpha_j^\vee \rangle
    = \frac{2(\omega_i, \alpha_j)}{(\alpha_j, \alpha_j)} = \delta_{ij},
    \quad 1 \leq i, j \leq \ell.
\]
They form a $\ZZ$-basis of the abstract weight lattice $\Lambda$, which sits between
the root lattice $\Lambda_r = \ZZ\Phi$ and $X(T)$ as follows:
\[
    \Lambda_r \subseteq X(T) \subseteq \Lambda,
\]
both inclusions being of finite index~\cite[Section~31.2]{humphreys_linear_1975}.
A weight $\chi = \sum_i c_i \omega_i \in \Lambda$ is dominant if and only if all
$c_i \geq 0$. In particular, when $G$ is simply connected one has $X(T) = \Lambda$,
and the dominant weights of $G$ are precisely the $\ZZ_{\geq 0}$-linear combinations
of $\omega_1, \ldots, \omega_\ell$.
\begin{exmp}[Dominant weights of $\SL{n}{K}$]
    For $G = \SL{n}{K}$ with the diagonal torus and the root system
    of \cref{ex:sln_roots}, the simple roots are $\alpha_i = \varepsilon_i - \varepsilon_{i+1}$.
    The fundamental dominant weights are
    \[
        \omega_k = \varepsilon_1 + \cdots + \varepsilon_k, \quad 1 \leq k \leq n-1.
    \]
    Using the restriction of the standard inner product on $\RR^n$, one computes
    $(\omega_k, \alpha_i) = \delta_{ik}$ and $(\alpha_i, \alpha_i) = 2$, giving
    $\langle \omega_k, \alpha_i^\vee \rangle = \delta_{ik}$ as required. Since
    $\SL{n}{K}$ is simply connected, $X(T) = \Lambda$ and the dominant weights
    are $\bigl\{\sum_k c_k \omega_k \mid c_k \in \ZZ_{\geq 0}\bigr\}$.
\end{exmp}

Finally, define a partial order on $X(T)$ by $\mu \leq \lambda$ if
$\lambda - \mu \in \ZZ_{\geq 0}\Delta$, i.e.\ $\lambda - \mu$ is a non-negative
integer combination of simple roots. This order underlies the notion of
``highest'' in the next section: in any $G$-module, a maximal vector will be seen
to carry the strictly largest weight for $\leq$, and the highest weight theorem will
then assert a bijection between irreducible rational $G$-modules and dominant weights.

\subsubsection{Maximal vector}

Let $\rho : G \rightarrow \GL{}{V}$ be a rational representation. By the Lie Kolchin theorem, $\rho(B)$ has a common eigenvector, so that there is a $1$-dimensional subspace of $V$ stable under $\rho(B)$. A vector spanning such a subspace is a \emph{maximal vector} of $V$. Clearly, it lies in some weight space $V_\lambda$ since $\rho(T)$ is contained in $\rho(B)$. The weights of the $G$-module generated by such a maximal vector are quite explicit.
\begin{theorem}
    Let $V$ be a $G$-module, $v^+$ a maximal vector of weight $\lambda$ and $V'$ the $G$-submodule generated by $v^+$. The weights of $V'$ are of the form $\lambda - \sum c_{\alpha} \alpha$, where $\alpha \in \Phi^+$ and $c_\alpha \in \ZZ_{\geq0}$.
\end{theorem}
A proof is found in~\cite[Section~31.2]{humphreys_linear_1975}.

The weight $\lambda$ of the maximal vector $v^+$ is called the \emph{highest weight} of $V'$: by the theorem, every other weight $\mu$ of $V'$ satisfies $\mu < \lambda$ in the partial order defined above (since $\lambda - \mu$ is a non-negative combination of positive roots). In particular, $V'$ admits no weight strictly larger than $\lambda$, justifying the name. Moreover, one may show that $\lambda$ is a dominant weight.

\subsubsection{The main theorem}

The above observations culminate in the following classification theorem, discussed in~\cite[Sections~31.2-31.4]{humphreys_linear_1975}.

\begin{theorem}[Highest weight theorem]
    \hfill
    \begin{enumerate}
        \item Every irreducible rational $G$-module $V$ has a unique (up to scalar) maximal vector, hence a unique highest weight $\lambda \in X(T)$, which is dominant.
        \item Two irreducible rational $G$-modules are isomorphic if and only if they have the same highest weight.
        \item For every dominant weight $\lambda \in X(T)$, there exists an irreducible rational $G$-module $L(\lambda)$ with highest weight $\lambda$.
    \end{enumerate}
\end{theorem}

The irreducible rational representations of $G$ are therefore in natural bijection with the set of dominant weights, parametrized by their non-negative integer coordinates with respect to the fundamental dominant weights $\omega_1, \ldots, \omega_l$.

While the explicit construction of these highest weight representations is quite involved, finding an irreducible representation with a specific highest weight guarantees that it is the only one with such a highest weight.

\begin{exmp}[Weights of $\Sym{k}{\CC^n}$ by the action of $\SLnC$]
    Let $G = \SLnC$, with $B$ the standard Borel subgroup, $T$ the maximal torus of diagonal matrices of determinant $1$. $G$ acts on $V = \Sym{k}{\CC^n}$ as described in \cref{subsubsection:actions}.\\
    Recall $X(T)$ is free abelian over $\varepsilon_1, \ldots, \varepsilon_{n-1}$, and $\varepsilon_n = - (\varepsilon_1 + \cdots + \varepsilon_{n-1})$. \\
    The weights of $\rho$ are found by analyzing the action of $T$ on a basis of $V$, such a basis being given by products $\bigodot_{i = 1}^n e_i^{\odot j_i}$ such that the $j_i$'s sum to $k$. Let $t = \textnormal{diag}(t_1, \ldots, t_n) \in T$. Then
    \[ t . \bigodot_{i = 1}^n e_i^{\odot j_i} = \bigodot_{i = 1}^n {(t.e_i)}^{\odot j_i} = \bigodot_{i = 1}^n {(t_i \cdot e_i)}^{\odot j_i} = \bigodot_{i = 1}^n {(\varepsilon_i(t) \cdot e_i)}^{\odot j_i} = \sum_{i=1}^{n}(j_i \cdot \varepsilon_i(t)) \cdot \bigodot_{i = 1}^n e_i^{\odot j_i}.\]   
    Thus the weights of $\rho$ are $\sum_{i = 1}^{n} j_i \varepsilon_i$, where the $j_i$'s sum to $n$. The ordering given by $B$ on the generators of $X(T)$ is $\varepsilon_1 > \cdots > \varepsilon_n$, which immediately implies that the highest weight of the representation of $\SLnC$ on $\Sym{k}{\CC^n}$ is $k \varepsilon_1$. In particular, we have shown that, for $\SLnC$ with standard Borel subgroup and maximal torus, $V(k \varepsilon_1) \cong \Sym{k}{\CC^n}$.
\end{exmp}
